\paragraph{One signed coefficient calculation.}
\label{j:local:atlas}\label{j:local:moments}
All bounds in this paragraph concern one actual marginal. Put
$a_\pm=(1\pm\beta)/2$, $b_\pm=(\rho\pm\gamma)/2$ and
$m=\E X_+=\E X_->0$. The moment matrices below give
$m^2\le a_+a_-$ and $b_\pm m\ge a_\pm^2$, hence
\[
 \frac{(1+\beta)^{3/2}}{\sqrt{1-\beta}}-\rho\le\gamma
 \le\rho-\frac{(1-\beta)^{3/2}}{\sqrt{1+\beta}},\qquad
 m\ge\max\left\{\frac{1+\beta^2}{2\rho},\frac{a_+^2}{b_+},\frac{a_-^2}{b_-}\right\}.
\]
The Gram matrix of $X,|X|,X|X|$ similarly gives
$\mu_4\ge\rho^2+(\gamma-\rho\beta)^2/(1-\beta^2)$.
Integrating $(x-a)^2(R-x)\ge0$ also gives
$2m\le(R-\rho+2a+Ra^2)/[a(2R+a)]$ for $a>0$; the cover uses
$a=2/10,\ldots,12/10$.

The same matrices give both signed constant deductions at once.
Write $\delta=\gamma-\rho\beta$ and
\[
 P_s=\frac{(\rho+\gamma)(\delta_+)^2}{8(1+\beta)^2}
       +\frac{(\rho-\gamma)((-\delta)_+)^2}{8(1-\beta)^2},\qquad
 P_g=\frac{\delta^4}{4\rho(1-\beta^2)^2}.
\]
Then $\rho\mu_4/2-\eta_5/4\le\rho^3/4-\max(P_s,P_g)$.
Indeed, on either sign, $c=\E X_s^4\ge b^2/a$ and
$f=\E X_s^5\ge c^2/b$, so
$\rho c/2-f/4\le\rho^2b/4-(c-\rho b)^2/(4b)$.
Distance to $c\ge b^2/a$ gives $P_s$; using
$\eta_5\ge\mu_4^2/\rho$ and the Gram bound gives $P_g$.
The maximum is a choice of two deductions, not their sum.
These moment statements allow an atom at zero. The following central-gap
clauses apply only after $(-A,B)$, $A,B>0$, has been acquired as below.

% Independent analytic license for the algebraic coefficient family.
% The surrounding text supplies the coefficient contract and existing notation.
\paragraph{Why the coefficient formulas bound one marginal.}
There are three underlying calculations. They use the same signwise
moments, conditional variances and gap budget throughout; the exact
coefficient object records their polynomial expansions.

First, on either sign let $Y=|X|\in[0,r]$ and write
$s_j=\E[Y^j;\text{that sign}]$. Both signs have $s_1=m$.
For $H_j=\left(\begin{smallmatrix}s_j&s_{j+1}\\s_{j+1}&s_{j+2}\end{smallmatrix}\right)$,
\[
 H_j\succeq0\quad(0\le j\le3),\qquad
 rH_j-H_{j+1}\succeq0\quad(0\le j\le2).
\]
Their quadratic forms integrate, respectively,
$Y^j(a+bY)^2$ and $Y^j(r-Y)(a+bY)^2$. These identities prove the
whole moment and support-contraction family, separately at the two
support radii. A quotient deduction from
$\left(\begin{smallmatrix}A&B\\B&C\end{smallmatrix}\right)\succeq0$
uses $C\ge B^2/A$ only when $A>0$. At $A=0$ it uses $B=0$, $C\ge0$;
an interval meeting zero may simply discard the nonnegative correction.
It may not divide by zero or square an unlicensed signed endpoint.

Second, the conditional Taylor identity
\eqref{j:local:re:conditional} integrates its polynomial remainder
against the nonnegative measures
$d\nu_\pm=(X-V)^2\ind_\pm dP$, of masses $\theta_\pm$.
With $a_\pm=|\E[X\mid\pm]|$, $\gamma_V=\E V^3$, it gives
\begin{equation}\label{j:local:entry201:Kexact}\begin{gathered}
 K_1(X)-K_1(V)=
 \int F(x,a_+;\rho+\gamma_V)\,d\nu_+
 -\int F(x,a_-;\rho-\gamma_V)\,d\nu_-,\\
 F(x,a;t)=x^2+2ax+3a^2-3-\tfrac32t(x+2a).
\end{gathered}\end{equation}
Indeed expand the product as
$\rho\gamma-\rho_V\gamma_V=\rho(\gamma-\gamma_V)
+\gamma_V(\rho-\rho_V)$ and use the conditional identities at orders
two, three and four. The factored derivatives certify $F_t\le0$ and
the positive Hessian $2(dx+da)^2+4(da)^2$.
Thus a rectangle maximum is at a corner; a minimum is among its corners,
its clipped edge stationary points and its interior stationary point.
All are included in the exact coefficient map. If the two signed
integrands lie in $[L,U]$, their sum lies in
$[\Theta\min(0,L),\Theta\max(0,U)]$ when
$\theta_++\theta_-\le\Theta$. This includes zero noise without division.

Third, direct integration over the central gap $(-A,B)$ gives
\begin{gather*}
 J_{1,c}=mP_1(A,B),\quad J_{2,c}=mP_2(A,B),\\
 P_1=B^3/3+AB^2/2-A^3/6,\qquad
 P_2=B^4/4+AB^3/3-A^4/12.
\end{gather*}
Both polynomials increase in $B$ and their $A$ derivatives change sign
only at $A=B$. Their extrema therefore use the indicated endpoints and
the clipped point $A=B$. For $G=\rho P_1-P_2$, the complete central
bound uses the degree-four tensor Bernstein coefficients on the closed
$(A,B)$ rectangle. The basis is nonnegative and sums to one, so retain
both $m_l$ and $m_u$ times every coefficient. This handles negative
coefficients and zero-width rectangles. The result is a maximum of
affine functions of $\rho$, hence convex. The checked inverse basis
identities prove coverage of the entire rectangle, not only a grid.

For a positive outer gap $(a,b)$, $a\ge B$, the square inequality
$(t-a)(\rho-t)\le(\rho-B)_+^2/4$ and signwise Tonelli give
\[
 3(\rho J_1-J_2)_o\le\tfrac38
 \{(\rho-B)_+^2T_++(R_--\rho)_+^2T_-\},\qquad
 -R_-T_-/2\le J_{1,o}\le R_+T_+/2.
\]
For the first bound on a negative gap write $x=-t$: its integrand is
$(v-x)x(x-\rho)$, whose positive part is at most
$x(R_--\rho)_+^2/4$. Enlarge each union of gaps to its complete
outer ray; the integrand is nonnegative, and Tonelli gives
$\int_B^\infty t\tau(t)\,dt=T_+/2$ and its negative counterpart.
This proves the factor $3/8$ even if the support has positive Lebesgue
measure. For the second bound use
$|t|\le R_\pm$ and $J_o^\pm\le T_\pm$.
Endpoint atoms do not change these Lebesgue integrals.

These costs share the one budget
$(\sigma-1)T_++(\sigma+1)T_-\le D-mQ$ from
\eqref{j:core:gap-defect}. Intersect it with any proved individual
caps before pricing a linear cost. Without caps, use
\[
 b_+T_++b_-T_-\le(D-mQ)
 \max\{0,b_+/(\sigma-1),b_-/(\sigma+1)\}.
\]
The zero option is necessary for signed costs. Sum costs before this
single pricing step; intersect alternative upper bounds afterwards.
Every coefficient interval is a necessary condition on this one law.
The analytic formation of the central gap and the prior support bounds
remain prerequisites; coefficient acceptance does not establish them.

\paragraph{The remaining endpoint recipes.}
For the conditional coefficient, use the mean and binary-cubic bounds
from $V=-m/q,m/p$ and $m^2=(1-\theta)pq$. Explicitly, if
$p\in[l,u]\subset(0,1/2)$ and $\theta\le\Theta$, the positive mean
lies between $\sqrt{(1-\Theta)(1-u)/u}$ and $\sqrt{(1-l)/l}$,
the negative magnitude mean between $\sqrt{(1-\Theta)l/(1-l)}$ and
$\sqrt{u/(1-u)}$, and
\[
 \sqrt{(1-\Theta)^3(1-2u)^2/[u(1-u)]}
 \le\gamma_V\le(1-2l)/\sqrt{l(1-l)}.
\]
In $F$, the edge stationary points are
$a=t/2-x/3$ and $x=3t/4-a$, clipped to their edges; the interior
point is $x=a=3t/8$. Retain the actual common $\rho$ interval.
Its lower bound may be raised to $\rho_B(u)$ only if
$3\min(A_l,B_l)\ge(3/2)\sup_I\rho_B$; the conditional cubic identity
and $(1-\theta)^{3/2}\ge1-3\theta/2$ then prove that license.
Intersect the signed $K_1$ interval with its previously proved error interval.

For the central constant, the cheaper alternative to the complete Bernstein
rectangle follows from
\[
 G_B=B(A+B)(\rho-B),\quad
 G(A,\rho;\rho)=\rho^4/12+\rho^3A/6-\rho A^3/6+A^4/12.
\]
The latter has derivatives $(\rho-A)^2(\rho+2A)/6$ in $A$ and
$\rho^2+\rho A$ twice in $\rho$, so use the upper $A$ endpoint and
both $\rho$ endpoints. Intersect the weighted outer budget with
$T_+\le b_{+,u}-m_lB_l^2$ and $T_-\le b_{-,u}-m_lA_l^2$.
For $U=J_1+(\mu_4-3)/6$, add the signed central and outer intervals
and the fourth-moment interval. Use the complete $K_1,U,\phi$ rectangle
in $|CK_1+\phi U|$; do not separate the two absolute values.

\begin{lemma}[Complete central rectangle]\label{j:local:entry201:central}
For $G=\rho P_1-P_2$, let $b_{ij}(\rho)$ be its degree-four tensor
Bernstein coefficients after mapping the proved $(A,B)$ rectangle to
$[0,1]^2$. Here one may take
\[
 A_u=\sqrt{u/(1-u)},\quad B_u=\sqrt{(1-l)/l},\quad
 m_l=\sqrt{(1-\Theta)\min_{p=l,u}pq},\quad
 m_u=\sqrt{\max_{p\in[l,u]}pq}.
\]
Then
\[
 \rho J_{1,c}-J_{2,c}\le
 \max_{0\le i,j\le4}\{m_l b_{ij}(\rho),m_u b_{ij}(\rho)\}.
\]
This upper bound is convex in $\rho$.
\end{lemma}
The preceding positive-basis argument proves the lemma, including closed
faces and degenerate rectangles. The fixed algebraic objects record the
seven moment matrices, the displayed derivative and square identities,
and all25 inverse basis identities. Exact coefficient comparison checks
these finite identities; the integral and common-budget arguments above
supply their mathematical meaning. No acquisition or Fourier target
follows from coefficient acceptance alone.
